\chapter*{ТЕРМИНЫ И ОПРЕДЕЛЕНИЯ}
\addcontentsline{toc}{chapter}{ТЕРМИНЫ И ОПРЕДЕЛЕНИЯ}

Видеоаналитика реального времени - совокупность методов и программных средств, обеспечивающих анализ видеопотока с ограниченной задержкой и выдачей результатов, пригодных для оперативного реагирования.

Детекция объекта - результат локализации и классификации объекта в кадре, содержащий ограничивающий прямоугольник, класс и оценку достоверности.

Сопровождение объекта - поддержание устойчивой траектории объекта между последовательными кадрами видеопотока.

Трек (траектория) - последовательность состояний одного и того же физического объекта в ряде кадров, объединенная единым устойчивым идентификатором.

Инференс модели - применение обученной модели к входному кадру для получения детекций или иных выходных признаков.

Событие - интерпретируемый факт, сформированный правиловым движком на основе траектории, пространственной геометрии сцены и временных условий.

Зона контроля - заданная область кадра, используемая для формализации ограничений и условий генерации событий.

Виртуальная линия - ориентированный отрезок (или ломаная) в координатах кадра, пересечение которого траекторией используется как триггер события.

Транзакционная таблица исходящих сообщений (transactional outbox) - архитектурный шаблон надежной доставки, при котором событие и запись о его последующей публикации фиксируются в одной транзакции в базе данных.

Идемпотентность доставки - свойство обработчика, при котором повторная обработка одного и того же сообщения не приводит к искажению итогового состояния внешней системы.

HTTP-получатель - внешняя система или тестовый сервис, принимающий сведения о событиях платформы по HTTP-запросам.

Очередь недоставленных сообщений - очередь брокера сообщений, в которую помещаются сообщения после исчерпания допустимого числа попыток доставки.
